%==========================================================================
% Section: Second application -- Wisconsin LGPIF
%   Results produced by lgpif_models.R; see results_lgpif/.
%==========================================================================
\section{Application to the Wisconsin Local Government Property Insurance Fund}
\label{sec.lgpif}

In this section, we apply Model~\ref{model.thin} to a second data set,
in which the units are insured entities rather than regions, the panel
is much larger, and the nested count is defined by the type of
coverage rather than by the claim size.

\subsection{Data}
\label{sec.lgpif.data}

The Wisconsin Local Government Property Insurance Fund (LGPIF) insures
local government entities, such as cities, counties, schools, towns, and
villages, against property losses. We use the data of E.~W.~Frees, which
record for each entity and year the claim counts of six coverage types:
building and contents (BC), inland marine (IM), and four types of
vehicle coverage, namely new and old vehicle collision (PN, PO) and new
and old vehicle comprehensive (CN, CO).
% TODO: add the reference for the data (e.g. Frees, Lee and Yang, 2016) and its source.
The training data cover \(5{,}677\) entity-years of \(1{,}234\) entities
over 2006--2010, and the hold-out data cover \(1{,}098\) entities in 2011,
of which all but \(9\) appear in the training data. We define
\[
\begin{aligned}
N_{it}&:=\text{number of claims over all six coverages},\\
M_{it}&:=\text{number of vehicle claims (PN, PO, CN, CO)},
\end{aligned}
\]
so that \(M_{it}\le N_{it}\) by construction, and \(M_{it}/N_{it}\) is the
share of vehicle claims among all claims of entity \(i\) in year \(t\).
Table~\ref{tab.lgpif.year} reports the annual totals. In contrast to the
application in Section~\ref{sec.empirical}, the share of vehicle claims
is stable over the years, between \(33\%\) and \(39\%\).

\begin{table}[htbp]
\centering
\caption{Annual totals of the LGPIF claim counts. The year 2011 is the
hold-out year.}
\label{tab.lgpif.year}
\small
\begin{tabular}{lrrrrrr}
\toprule
Year & 2006 & 2007 & 2008 & 2009 & 2010 & 2011\\
\midrule
Entities & 1{,}162 & 1{,}147 & 1{,}134 & 1{,}117 & 1{,}117 & 1{,}098\\
\(\sum_i N_{it}\) & 1{,}415 & 1{,}676 & 1{,}471 & 1{,}675 & 1{,}962 & 1{,}733\\
\(\sum_i M_{it}\) & 518 & 554 & 566 & 569 & 682 & 612\\
Share (\%) & 36.6 & 33.1 & 38.5 & 34.0 & 34.8 & 35.3\\
\bottomrule
\end{tabular}
\end{table}

The covariate vector \(x_{it}\) consists of an intercept, indicators of
the entity type with cities as the baseline, the indicator IsRC
provided in the data, the logarithm of the BC deductible, and
\(\log(1+c)\) for the BC, IM, and total vehicle coverage amounts \(c\), in
millions. We do not use the no-claim credit variables in the data,
because they are themselves based on the claim experience of the entity,
which is accounted for by the random effects.
% TODO: state the definition of IsRC.

Most entities have no vehicle coverage: \(62\%\) of the entity-years have
zero vehicle coverage amount, so that the probability of a vehicle claim
is zero by construction. We therefore apply the binomial part of
Model~\ref{model.thin} only to the entity-years with vehicle coverage,
and set the predicted vehicle claim count to zero for the entities
without it. Only \(4\) entity-years in the training data, with \(9\)
claims in total, record vehicle claims without vehicle coverage; they
are excluded from the binomial part. The binomial part is therefore
estimated from \(1{,}110\) entity-years of \(391\) entities with vehicle
coverage and at least one claim, which contain \(6{,}294\) claims, of
which \(2{,}880\) are vehicle claims.

\subsection{Estimation}
\label{sec.lgpif.estimation}

The estimation follows Section~\ref{sec.emp.estimation}, with the entity
in place of the region. For the binomial part, the Gauss--Hermite
quadrature uses \(200\) nodes. Since some entities have several hundred
claims over the training period, their likelihood contributions are
sharply peaked in \(\theta\); with \(200\) nodes, \(\ell^{(2)}\) agrees
with that computed by \(400\) nodes and by adaptive numerical
integration to within \(5\times10^{-5}\).

Table~\ref{tab.lgpif.coef} reports the estimates. Both random effects
are strongly supported. For the total claim count, the variance of the
gamma random effect is \(1/\widehat r=0.735\), the Pearson statistic of
the Poisson GLM is \(5.18\) times its degrees of freedom, and the
log-likelihood increases by \(3{,}176\). For the vehicle claims, the
variance of the Gaussian random effect is \(\widehat\sigma^2=0.360\),
the Pearson statistic of the probit GLM is \(2.15\) times its degrees of
freedom, and the log-likelihood increases by \(395\). Both likelihood
ratio tests have \(p\)-values below \(10^{-100}\). Several coefficients
change markedly once the random effects are introduced, and their GLM
standard errors are considerably smaller than the GLMM ones, which
reflects the neglected dependence among the repeated observations of an
entity in the GLMs.

\begin{table}[htbp]
\centering
\caption{Estimates for the LGPIF data, with standard errors in
parentheses. Total claims: Poisson GLM and Poisson-gamma GLMM; vehicle
claims: probit GLM and probit GLMM.}
\label{tab.lgpif.coef}
\small
\begin{tabular}{lrrrr}
\toprule
& \multicolumn{2}{c}{Total claims \(N_{it}\)} & \multicolumn{2}{c}{Vehicle claims \(M_{it}\)}\\
\cmidrule(lr){2-3}\cmidrule(lr){4-5}
& GLM & GLMM & GLM & GLMM\\
\midrule
Intercept & \(-1.611\) (0.065) & \(-1.093\) (0.179) & \(1.090\) (0.106) & \(-0.548\) (0.261)\\
TypeCounty & \(0.791\) (0.037) & \(0.478\) (0.170) & \(0.915\) (0.066) & \(0.831\) (0.172)\\
TypeMisc & \(-1.463\) (0.097) & \(-0.648\) (0.167) & \(0.236\) (0.277) & \(0.152\) (0.369)\\
TypeSchool & \(-0.655\) (0.050) & \(-0.486\) (0.128) & \(0.496\) (0.105) & \(0.405\) (0.204)\\
TypeTown & \(-0.225\) (0.094) & \(-0.425\) (0.156) & \(-0.998\) (0.176) & \(-0.610\) (0.267)\\
TypeVillage & \(0.006\) (0.049) & \(0.014\) (0.115) & \(-0.263\) (0.095) & \(-0.097\) (0.176)\\
IsRC & \(-0.343\) (0.025) & \(-0.036\) (0.073) & \(0.011\) (0.037) & \(-0.133\) (0.089)\\
\(\log(1+\text{BC coverage})\) & \(0.856\) (0.017) & \(0.702\) (0.040) & \(-0.536\) (0.028) & \(-0.350\) (0.062)\\
log(BC deductible) & \(-0.121\) (0.009) & \(-0.191\) (0.024) & \(-0.005\) (0.014) & \(0.137\) (0.034)\\
\(\log(1+\text{IM coverage})\) & \(-0.399\) (0.031) & \(0.198\) (0.096) & \(0.017\) (0.052) & \(-0.068\) (0.142)\\
\(\log(1+\text{vehicle coverage})\) & \(0.516\) (0.015) & \(0.487\) (0.052) & \(0.350\) (0.041) & \(0.377\) (0.086)\\
\(1/r\) or \(\sigma^2\) & -- & \(0.735\) (0.048) & -- & \(0.360\) (0.053)\\
\midrule
Observations & \multicolumn{2}{c}{\(5{,}677\)} & \multicolumn{2}{c}{\(1{,}110\)}\\
Log-likelihood & \(-9{,}160.5\) & \(-5{,}984.2\) & \(-1{,}700.8\) & \(-1{,}305.3\)\\
AIC & \(18{,}342.9\) & \(11{,}992.3\) & \(3{,}423.6\) & \(2{,}634.7\)\\
BIC & \(18{,}416.0\) & \(12{,}072.1\) & \(3{,}478.7\) & \(2{,}694.8\)\\
Pearson \(\chi^2\)/df & \(5.18\) & -- & \(2.15\) & --\\
\bottomrule
\end{tabular}
\end{table}

The estimates are consistent with the coverage structure: a larger
vehicle coverage is associated with more claims in total and with a
higher share of vehicle claims, whereas a larger BC coverage is
associated with more claims in total but a lower share of vehicle
claims, since it increases the number of BC claims. Counties have both
more claims and a higher share of vehicle claims than cities, and towns
have a lower share. As in Section~\ref{sec.empirical}, these
associations are descriptive.

\subsection{One-step-ahead prediction of 2011}
\label{sec.lgpif.prediction}

We predict the claim counts of 2011 from the training history as in
Section~\ref{sec.emp.prediction}: the total claim count by
\(\E{N_{i,2011}\mid\mathcal H_{i,2010}}\), and the vehicle claim count
(a) given the realized total claim count by
\(N_{i,2011}\,\widehat p_{i,2011}\), and (b) without it by
\(\E{N_{i,2011}\mid\mathcal H_{i,2010}}\,\widehat p_{i,2011}\). The
\(9\) entities without training data receive the a priori predictions.
Table~\ref{tab.lgpif.crit} reports the results.

\begin{table}[htbp]
\centering
\caption{Prediction criteria for 2011 (\(1{,}098\) entities, of which
\(405\) have vehicle coverage). The binomial deviance is computed over the
entities with vehicle coverage and \(N_{i,2011}\ge1\), and the Poisson
deviance of \(M\) over the entities with vehicle coverage.}
\label{tab.lgpif.crit}
\small
\begin{tabular}{llrr}
\toprule
& Criterion & GLM & GLMM\\
\midrule
Total claims \(N\) & RMSE & \(7.401\) & \(3.563\)\\
& MAE & \(1.473\) & \(0.941\)\\
& Poisson deviance & \(2{,}902.7\) & \(1{,}230.8\)\\
& Predicted total (actual \(1{,}733\)) & \(1{,}780.2\) & \(1{,}880.2\)\\
\midrule
Vehicle claims \(M\), (a) given \(N_{i,2011}\) & RMSE & \(0.898\) & \(0.488\)\\
& MAE & \(0.212\) & \(0.142\)\\
& Binomial deviance & \(426.0\) & \(270.4\)\\
& Predicted total (actual \(612\)) & \(627.3\) & \(646.0\)\\
\midrule
Vehicle claims \(M\), (b) unconditional & RMSE & \(2.137\) & \(1.208\)\\
& MAE & \(0.490\) & \(0.290\)\\
& Poisson deviance & \(792.2\) & \(343.2\)\\
& Predicted total (actual \(612\)) & \(598.2\) & \(693.2\)\\
\bottomrule
\end{tabular}
\end{table}

In this application, the GLMM improves all entity-level criteria
substantially. For the vehicle claims given \(N_{i,2011}\), the RMSE
decreases from \(0.898\) to \(0.488\) and the binomial deviance from
\(426.0\) to \(270.4\); without \(N_{i,2011}\), the RMSE decreases from
\(2.137\) to \(1.208\). The credibility adjustments are large: the
posterior means \(\E{\Theta_i^{(1)}\mid\mathcal H_{i,2010}}\) range from
\(0.05\) to \(27.3\) with median \(0.80\), and the posterior means of
\(\Theta_i^{(2)}\) range from \(-1.92\) to \(1.31\) among the entities with
vehicle coverage. The improvement is not driven by a few large
entities. Excluding the \(11\) entities in the top \(1\%\) of
\(N_{i,2011}\), the RMSEs of the GLM and the GLMM are \(2.57\) and
\(1.89\) for \(N\), and \(0.92\) and \(0.65\) for \(M\) given
\(N_{i,2011}\). In \(2{,}000\) bootstrap resamples of the entities, the
GLMM has the smaller RMSE in \(97.0\%\) of the resamples for \(N\) and in
all resamples for \(M\), in both cases (a) and (b). The GLMMs
over-predict the portfolio totals, by \(8.5\%\) for \(N\) and \(13.3\%\)
for \(M\) in case (b), whereas the GLMs predict the totals more closely;
the entity-level accuracy, which matters for experience rating, is
nevertheless much higher for the GLMMs.

\subsection{Comparison of the two applications}
\label{sec.lgpif.comparison}

The two applications lead to different conclusions about the benefit of
the experience rating of the nested count. In the LGPIF data, the share
of vehicle claims of an entity is a persistent characteristic,
determined largely by the composition of its insured property, and the
annual share in the portfolio is stable. The static random effect
\(\Theta_i^{(2)}\) captures this persistence, and the credibility
predictions of Theorem~\ref{thm.thin.pred} improve on the GLM
substantially. In the regional data of Section~\ref{sec.empirical}, the
share of extreme claims is dominated by year-specific shocks that are
common to the regions, and the experience of a region is only weakly
informative about its next-year share. The static model is therefore
well suited to nested counts whose split reflects stable unit
characteristics, while nested counts driven by common shocks call for
the extensions discussed in Section~\ref{sec.emp.discussion}.
